\documentclass[10pt,a4paper]{amsart}
\usepackage[margin=19mm]{geometry}
\usepackage{amsmath,amssymb,mathtools}
\usepackage[scr=rsfs,cal=euler]{mathalfa}
\usepackage{xfrac}
\usepackage[unicode,hidelinks,bookmarks=false]{hyperref}
\newcommand{\ie}{i.e.\ }
\newcommand{\cf}{cf.\ }
\newcommand{\resp}{respectively\ }
\newcommand{\Subsection}{\subsection}
\providecommand{\nobreakdash}{\nobreak\hbox{-}}
\setlength{\emergencystretch}{2em}
\multlinegap0pt
% ---------------------------------------------------------------------------
% Editorial AMS wrapper prelude. The theorem environments and the mathematical macros below are
% transcribed from the paper's own preamble (cached-originals/source0.tex lines 44--113); the AMS amsart
% class, the named format packages of the pinned TeX Live runtime and this editorial formatting are not
% part of the author's text. No mathematics is changed.
% ---------------------------------------------------------------------------
\newcommand{\mn}[1]{\marginpar{\tiny #1}}
\def\red{\color{red}}
\def\green{\color{green}}
\def\blue{\color{blue}}
\def\cyan{\color{cyan}}
\def\magenta{\color{magenta}}
\def\yellow{\color{yellow}}


\DeclareMathOperator{\Td}{Td}
\DeclareMathOperator{\ch}{ch}
\DeclareMathOperator{\cc}{c}
\DeclareMathOperator{\Tr}{Tr}
\DeclareMathOperator{\ric}{Ric}
\DeclareMathOperator{\End}{End}



\newtheorem{maindefn}{{\sc Definition}}
\newtheorem*{maindefn*}{{\sc Definition}}
\newtheorem{theo}{{\sc Theorem}}[section]
\newtheorem{mainlem}{{\sc Lemma}}
\newtheorem*{maintheo}{{\sc Theorem}}
\newtheorem{conj}{{\sc Conjecture}}
\newtheorem{mainprop}{{\sc Proposition}}
\newtheorem{cor}[theo]{{\sc Corollary}}
\newtheorem{prop}[theo]{{\sc Proposition}}
\newenvironment{example}{\medskip\noindent{\it Example:\/} }{\medskip}

\newtheorem{mainthm}{Theorem}
\renewcommand{\themainthm}{\Alph{mainthm}}


\theoremstyle{definition}
\newtheorem{definition}{Definition}

\newtheorem{theorem}{Theorem}
\newtheorem{lemma}{Lemma}
\newtheorem{proposition}{Proposition}


\theoremstyle{remark}
\newtheorem{remark}{Remark}

\newenvironment{Ack}{\medskip\noindent{\it Acknowledgements:\/} }{\medskip}
\newtheorem*{mycor}{Corollary}


\newenvironment{claim}{\medskip\noindent{\it Claim:\/} }{\medskip}



\begin{document}
\title{Asymptotic expansion of the variation of the Quillen metric and its moment map interpretation}
\author{Kiyoon Eum}
\date{Source: arXiv:2510.25456v3; distributed under CC BY 4.0}
\maketitle
\noindent\emph{Source, version and licence.} \emph{Asymptotic expansion of the variation of the Quillen metric and its moment map interpretation}, by Kiyoon Eum. The source of this editable unit is the e-print of \textbf{version v3} of arXiv:2510.25456, https://arxiv.org/abs/2510.25456v3, whose material is distributed under the Creative Commons Attribution 4.0 International licence, http://creativecommons.org/licenses/by/4.0/, as recorded in the arXiv licence metadata for this paper and shown on that abstract page. The whole of the body below is version v3, the newest version of the paper. Full rights notice: licenses/arxiv-2510.25456-CC-BY-4.0.txt.\par\smallskip
\noindent\textit{Editorial scope.} Counted unit: original Proposition 1 of the article, the statement carried by the source environment \texttt{proposition}, cached-originals/source0.tex lines 457--459, printed as \textbf{Proposition 1} exactly as in the article. It is delivered together with the author's own complete proof as the single primary proof body, source0.tex lines 461--475 (15 lines): the \texttt{proof} environment that immediately follows the statement and proves the equivariance of the determinant line bundle and the invariance of the Quillen metric. No proof marker is added and no mathematics is written, repaired or re-derived; the author's notation and the printed slips of the source are kept literal. Necessary complete context is retained uncounted, in source order: source0.tex lines 403--417, the definition of the functionals $\widetilde{Z}_j(M,\omega)$ whose variation the article studies; and lines 435--440, the definition of the coefficients $b_{j-1}$ and of the exponent $\widetilde{t}_j$ used in the same formulas. The article's own numbering is reproduced by explicit counter settings: this article declares each theorem-like environment with its own global counter (\texttt{\textbackslash newtheorem\{proposition\}\{Proposition\}} and likewise for theorem, definition, lemma and remark), and every one of those counters is set before each retained passage, so the counted statement prints as Proposition 1 exactly as in the article. No \texttt{\textbackslash cite} command occurs in any retained passage and no \texttt{\textbackslash ref} or \texttt{\textbackslash eqref} inside a retained passage points outside it: the closure of the retained passages over references and citations was computed, it contains no reference and no citation at all, so this unit delivers no bibliography entry and the article's own bibliography data is neither spliced into the bundled source nor redistributed. The retained passages contain no figure environment and no graphics-inclusion command, so no artwork is redistributed. The source's own class is AMS \texttt{amsart} (\texttt{\textbackslash documentclass[12pt,reqno,a4paper]\{amsart\}}); the e-print ships no class or style file, so no publisher class is shipped, used or redistributed, and the wrapper is the same AMS \texttt{amsart} class with the paper's own macros and theorem declarations transcribed from its preamble. Every declared body of this unit is an exact, unmodified span of source0.tex: no body is a bounded passage comparison, \texttt{omit} was not used anywhere, and consequently no fidelity receipt is asserted in delivery-evidence.json. The complete concatenated original source is bundled as sources/188037/source0.tex. Omitted from this editable unit and therefore absent from it are source0.tex lines 1--402; 418--434; 441--456; 460--460; 476--654, that is the article's own preamble, title block and abstract, the introduction with the two main theorems, the sections on the space of complex structures and on the moment map, the appendix, and the article's bibliography data, all of which remain present in the bundled source but are not delivered here. The AMS \texttt{amsart} wrapper, the named format packages of the pinned TeX Live runtime and this editorial source, licence and scope paragraph are editorial formatting only.\par\smallskip
\setcounter{section}{3}
\setcounter{equation}{4}
\setcounter{conj}{0}
\setcounter{definition}{1}
\setcounter{lemma}{0}
\setcounter{maindefn}{0}
\setcounter{mainlem}{0}
\setcounter{mainprop}{0}
\setcounter{mainthm}{2}
\setcounter{proposition}{0}
\setcounter{remark}{1}
\setcounter{theo}{0}
\setcounter{theorem}{2}
\begin{definition}
We define
\begin{equation*}
\widetilde{Z}_j(M,\omega) := \widetilde{t}_j+\widetilde{\ell}_j\in C^\infty(M).
\end{equation*}
Then we have
\begin{align}\label{quilvar}
&\frac{\delta}{\delta\varphi} \int_M iBC(\Td_j;g_{\varphi},g)\frac{\omega^{n+1-j}}{(n+1-j)!}+\Td_{j}(R_\varphi)\frac{1}{(n+1-j)!}\sum^{n-j}_{s=0}\varphi \omega_{\varphi}^s\wedge \omega^{n-j-s}\nonumber\\
&=\frac{1}{(n-j)!}\int_M \delta\varphi \,\widetilde{Z}_j(M,\omega)\, \omega^n
\end{align}
and
\begin{equation*}
\int_M \widetilde{Z}_j(M,\omega) \omega^n = \int_M \Td_j(M) [\omega]^{n-j}.
\end{equation*}
\end{definition}

\setcounter{section}{3}
\setcounter{equation}{7}
\setcounter{conj}{0}
\setcounter{definition}{2}
\setcounter{lemma}{0}
\setcounter{maindefn}{0}
\setcounter{mainlem}{0}
\setcounter{mainprop}{0}
\setcounter{mainthm}{2}
\setcounter{proposition}{0}
\setcounter{remark}{1}
\setcounter{theo}{0}
\setcounter{theorem}{2}
\begin{definition}
Define $b_{j-1}$ for each $j\geq 1$ by
\begin{equation*}
\frac{n!}{(n-j)!}\widetilde{Z}_j(M,\omega)=a_j(g)-\Delta a_{j-1}(g)+b_{j-1}(g).
\end{equation*}
\end{definition}

\setcounter{section}{3}
\setcounter{equation}{8}
\setcounter{conj}{0}
\setcounter{definition}{3}
\setcounter{lemma}{0}
\setcounter{maindefn}{0}
\setcounter{mainlem}{0}
\setcounter{mainprop}{0}
\setcounter{mainthm}{2}
\setcounter{proposition}{0}
\setcounter{remark}{1}
\setcounter{theo}{0}
\setcounter{theorem}{2}
\begin{proposition}
The determinant line bundle $\lambda^{(k)}\rightarrow \mathcal{J}_{int}$ is $\mathcal{G}$-equivariant and $\mathcal{G}$ acts on $\lambda^{(k)}$ by biholomorphisms preserving the Quillen metric $||\cdot||_{\omega,\widetilde{h}^k}$ constructed from the data $g^M=g_\omega, h^\xi=\widetilde{h}^k$.
\end{proposition}

\setcounter{section}{3}
\setcounter{equation}{8}
\setcounter{conj}{0}
\setcounter{definition}{3}
\setcounter{lemma}{0}
\setcounter{maindefn}{0}
\setcounter{mainlem}{0}
\setcounter{mainprop}{0}
\setcounter{mainthm}{2}
\setcounter{proposition}{1}
\setcounter{remark}{1}
\setcounter{theo}{0}
\setcounter{theorem}{2}
\begin{proof}
Let $\psi\in\mathcal{G}$. By the definition of $\mathcal{G}$, the action of $\psi$ on $\Gamma(M,L^k)$ intertwines $D^2_J$ and $D^2_{\psi.J}$ and maps the eigenspace of $D^2_J$ onto the eigenspace of $D^2_{\psi.J}$ with the same eigenvalue. Hence, the line bundle $\lambda^{(k)} \rightarrow \mathcal{J}_{int}$ is $\mathcal{G}$-equivariant. 

To show that $\mathcal{G}$ acts on $\lambda^{(k)}$ by biholomorphisms, we prove that $\mathcal{G}$ preserves $\widetilde{\nabla}$ on $E^p$ and $P^b:E\rightarrow K^b$ on $U^b$. By the definition of $\widetilde{L}^k\rightarrow\mathcal{X}$, for a smooth section $e$ of $E^p$ and $A\in T\mathcal{J}_{int}$,
\begin{equation*}
\widetilde{\nabla}_A e(J) = \nabla_{A^H} e(z;J)=\nabla_{0\oplus A}e(z;J)=\Pi_J^{0,p}\left.\frac{d}{dt}\right|_{t=0} e(z;J+tA)=:\delta_A e(J),
\end{equation*}
where we used the notation $z\in M, J\in \mathcal{J}_{int}$ to indicate the dependence of sections on $M,\mathcal{J}_{int}$ and $\Pi_J^{0,p}$ to denote the projection onto the $(0,p)$-part. This implies
\begin{equation*}
\widetilde{\nabla}_A (\psi.e)=\delta_A \psi.(e(\psi^{-1}.z;\psi^{-1}.J))=\psi. \delta_{\psi^{-1}_*A}e=\psi. \widetilde{\nabla}_{\psi^{-1}_*A}e.
\end{equation*}
Thus $\mathcal{G}$ preserves $\widetilde{\nabla}$. We have already shown that $\psi$ maps $K^b_J$ onto $K^b_{\psi.J}$. Again, by the definition of $\mathcal{G}$, $\psi$ preserves the $L^2$ Hermitian product on $E$. Hence, $\psi$ preserves the orthogonal projection operators $P^b$ over $U^b$. 

This proves that ${}^0\nabla^b$ is $\mathcal{G}$-invariant over $U^b$. Since $\mathcal{G}$ acts on $\mathcal{J}_{int}$ by biholomorphisms, this implies that $\mathcal{G}$ acts on $\lambda^{(k)}$ by biholomorphisms. Finally, the spectra of $D^2_J$ and $D^2_{\psi.J}$ are identified by the action of $\psi$, and hence the Ray--Singer analytic torsion is $\mathcal{G}$-invariant. This proves the $\mathcal{G}$-invariance of the Quillen metric.
\end{proof}

\end{document}
